\chapter{Akar Kuadrat}\label{ch:g9:sqrt}

Berapa panjang diagonal persegi bersisi $1$? Teorema Pythagoras
menjawab $\sqrt 2$, yaitu bilangan yang kuadratnya $2$ --- dan yang
ternyata bukan \omterm{def:g9:fractions:fraction}{pecahan}. Bab ini mendefinisikan \omterm{def:g9:sqrt:def}{akar kuadrat},
menetapkan aturan berhitung dengannya, lalu mengajarkan cara
menyederhanakan bentuk seperti $\sqrt{75}$.

\section{Definisi dan sifat pertamanya}

\begin{definition}[Akar kuadrat]\label{def:g9:sqrt:def}
Misalkan $a \geq 0$. \emph{Akar kuadrat}\index{akar kuadrat} dari $a$,
yang ditulis $\sqrt a$, adalah satu-satunya bilangan
\emph{taknegatif} yang kuadratnya $a$:
\[
\sqrt a \geq 0
\qquad\text{dan}\qquad
\left(\sqrt a\right)^2 = a .
\]
Bilangan negatif tidak punya akar kuadrat, karena setiap kuadrat bernilai taknegatif.
\end{definition}

\begin{example}
$\sqrt{49} = 7$, $\sqrt 0 = 0$, $\sqrt 1 = 1$,
$\sqrt{2.25} = 1.5$. \omterm{def:g9:sqrt:def}{Akar kuadrat} pertama yang perlu dihafal adalah
akar \emph{kuadrat sempurna}: $1, 4, 9, 16, 25, 36, 49, 64, 81, 100,
121, 144$.
\end{example}

\begin{omfigure}
\begin{tikzpicture}[scale=1.35]
  \draw[thick, fill=omDef!8] (0,0) rectangle (2,2);
  \draw[omThm, thick] (0,0) -- (2,2);
  \node[below, font=\small] at (1,0) {$1$};
  \node[left, font=\small] at (0,1) {$1$};
  \node[omThm, above left=-2pt, font=\small, rotate=45] at (1.08,0.92)
    {$\sqrt2$};
  \draw[thin] (1.75,0) -- (1.75,0.25) -- (2,0.25);
\end{tikzpicture}

{\small Diagonal persegi satuan punya panjang $\sqrt{1^2 + 1^2} =
\sqrt2 \approx 1.414$, menurut teorema Pythagoras
(\cref{ch:g9:trig} mengingatkannya). Bilangan ini irasional:
desimalnya tidak pernah berulang.}
\end{omfigure}

\begin{proposition}[Akar kuadrat sebuah kuadrat]\label{prop:g9:sqrt:square}
Untuk setiap bilangan real $x$ (positif atau bukan):
\[
\sqrt{x^2} = \abs{x} =
\begin{cases}
x & \text{kalau } x \geq 0,\\
-x & \text{kalau } x < 0 .
\end{cases}
\]
\end{proposition}

\begin{proof}
Bilangan $\abs{x}$ bernilai taknegatif dan kuadratnya $x^2$ (sebuah
bilangan dan \omterm{def:g7:negatives:def}{lawannya} punya kuadrat yang sama). Karena ia
satu-satunya bilangan taknegatif yang kuadratnya $x^2$, maka ia
$\sqrt{x^2}$.
\end{proof}

\begin{example}
$\sqrt{(-5)^2} = \sqrt{25} = 5 = \abs{-5}$: \omterm{def:g9:sqrt:def}{akar kuadratnya}
``melupakan'' tandanya, bukan memulihkannya. Khususnya
$\sqrt{x^2} = x$ itu \emph{keliru} untuk $x$ yang negatif.
\end{example}

\section{Hasil kali dan hasil bagi akar kuadrat}

\begin{theorem}[Aturan perkalian dan pembagian]\label{thm:g9:sqrt:rules}
Untuk semua $a \geq 0$ dan $b \geq 0$:
\[
\sqrt{a b} = \sqrt a \times \sqrt b,
\qquad\text{dan untuk } b > 0:\quad
\sqrt{\frac ab} = \frac{\sqrt a}{\sqrt b} .
\]
\end{theorem}

\begin{proof}
Bilangan $\sqrt a \times \sqrt b$ bernilai taknegatif (yaitu \omterm{def:g2:mult:def}{hasil
kali} bilangan taknegatif), dan kuadratnya adalah
\[
\left(\sqrt a \times \sqrt b\right)^2
= \left(\sqrt a\right)^2 \times \left(\sqrt b\right)^2
= a b .
\]
Karena ketunggalan bilangan taknegatif yang kuadratnya $ab$, bilangan
itu sama dengan $\sqrt{ab}$. Aturan \omterm{def:g3:division:remainder}{hasil baginya} dibuktikan dengan
cara yang sama.
\end{proof}

\begin{remark}[Tidak ada aturan seperti itu untuk jumlah!]
$\sqrt{a + b}$ \emph{bukan} $\sqrt a + \sqrt b$ pada umumnya:
\[
\sqrt{9 + 16} = \sqrt{25} = 5,
\qquad\text{tetapi}\qquad
\sqrt 9 + \sqrt{16} = 3 + 4 = 7 .
\]
\end{remark}

\begin{method}[Menyederhanakan $\sqrt n$]\label{met:g9:sqrt:simplify}
Untuk menyederhanakan \omterm{def:g9:sqrt:def}{akar kuadrat} sebuah bilangan bulat:
\begin{enumerate}
  \item carilah kuadrat sempurna terbesar yang membagi $n$ (faktorkan
        $n$ kalau perlu);
  \item pisahkan dengan aturan \omterm{def:g2:mult:def}{hasil kalinya} lalu keluarkan
        kuadratnya: $\sqrt{k^2 m} = k\sqrt m$;
  \item periksalah bahwa tidak ada kuadrat sempurna yang membagi apa
        yang tertinggal di bawah akarnya.
\end{enumerate}
\end{method}

\begin{example}\label{ex:g9:sqrt:simplify}
Sederhanakan $\sqrt{75}$: karena $75 = 25 \times 3$,
\[
\sqrt{75} = \sqrt{25 \times 3} = \sqrt{25} \times \sqrt 3 = 5\sqrt3 .
\]
Sederhanakan $\sqrt{72}$: karena $72 = 36 \times 2$, berlaku
$\sqrt{72} = 6\sqrt2$. Dan sebuah \omterm{def:g3:division:remainder}{hasil bagi}:
$\sqrt{\dfrac{49}{16}} = \dfrac{\sqrt{49}}{\sqrt{16}} = \dfrac74$.
\end{example}

\begin{example}[Menjumlahkan akar kuadrat]\label{ex:g9:sqrt:add}
\omterm{def:g1:addition:def}{Jumlah} \omterm{def:g9:sqrt:def}{akar kuadrat} menyederhana hanya ketika akarnya \emph{sejenis}.
Hitunglah $\sqrt{12} + \sqrt{27} - \sqrt{48}$, dengan menyederhanakan
tiap sukunya lebih dahulu:
\begin{align*}
\sqrt{12} &= \sqrt{4 \times 3} = 2\sqrt3, \\
\sqrt{27} &= \sqrt{9 \times 3} = 3\sqrt3, \\
\sqrt{48} &= \sqrt{16 \times 3} = 4\sqrt3,
\end{align*}
jadi \omterm{def:g1:addition:def}{jumlahnya} adalah $2\sqrt3 + 3\sqrt3 - 4\sqrt3 = (2 + 3 - 4)\sqrt3 = \sqrt3$.
\end{example}

\begin{example}[Menjabarkan dengan akar kuadrat]\label{ex:g9:sqrt:expand}
Kesamaan aljabar berlaku pula untuk \omterm{def:g9:sqrt:def}{akar kuadrat}. Jabarkan
$\left(\sqrt5 + 2\right)^2$ dengan $(a+b)^2 = a^2 + 2ab + b^2$:
\[
\left(\sqrt5 + 2\right)^2 = 5 + 2 \times 2\sqrt5 + 4 = 9 + 4\sqrt5 .
\]
Dan dengan kesamaan yang ketiga, $(a+b)(a-b) = a^2 - b^2$:
\[
\left(\sqrt7 + \sqrt3\right)\left(\sqrt7 - \sqrt3\right) = 7 - 3 = 4 :
\]
\omterm{def:g2:mult:def}{hasil kali} dua bilangan irasional dapat menjadi bilangan bulat.
\end{example}

\section{Persamaan \texorpdfstring{$x^2 = a$}{x2 = a}}

\begin{theorem}[Menyelesaikan $x^2 = a$]\label{thm:g9:sqrt:equation}
\begin{enumerate}
  \item Kalau $a > 0$, \omterm{def:g8:equations:def}{persamaan} $x^2 = a$ punya persis dua
        penyelesaian: $\sqrt a$ dan $-\sqrt a$;
  \item kalau $a = 0$, satu-satunya penyelesaiannya adalah $0$;
  \item kalau $a < 0$, tidak ada penyelesaiannya.
\end{enumerate}
\end{theorem}

\begin{proof}
Tulislah ulang $x^2 = a$ menjadi $x^2 - \left(\sqrt a\right)^2 = 0$
ketika $a \geq 0$, lalu faktorkan sebagai \omterm{ex:g1:subtraction:difference}{selisih} kuadrat:
\[
\left(x - \sqrt a\right)\left(x + \sqrt a\right) = 0 ,
\]
jadi $x = \sqrt a$ atau $x = -\sqrt a$ (keduanya berimpit ketika
$a = 0$). Ketika $a < 0$ tidak ada penyelesaiannya, karena
$x^2 \geq 0 > a$ untuk setiap $x$.
\end{proof}

\begin{omfigure}
\begin{tikzpicture}
\begin{axis}[omaxis,
  width=7.6cm, height=5.4cm,
  xmin=-3.1, xmax=3.1, ymin=-1.6, ymax=5.4,
  xtick={-1.732,1.732}, xticklabels={$-\sqrt3$,$\sqrt3$},
  ytick={3}, yticklabels={$3$}]
  \addplot[omDef, thick, domain=-2.25:2.25, samples=90] {x^2};
  \addplot[omThm, thick, domain=-2.9:2.9] {3};
  \fill (axis cs:-1.732,3) circle (1.5pt);
  \fill (axis cs:1.732,3) circle (1.5pt);
  \draw[dashed] (axis cs:-1.732,3) -- (axis cs:-1.732,0);
  \draw[dashed] (axis cs:1.732,3) -- (axis cs:1.732,0);
\end{axis}
\end{tikzpicture}

{\small Menyelesaikan $x^2 = 3$ pada parabola $y = x^2$: garis mendatar
$y = 3$ memotongnya di kedua absis $\pm\sqrt3$.}
\end{omfigure}

\begin{example}
$x^2 = 16$: penyelesaiannya $4$ dan $-4$. \quad
$x^2 = 5$: penyelesaiannya $\sqrt5$ dan $-\sqrt5$ (nilai persisnya;
$\sqrt5 \approx 2.236$). \quad
$x^2 = -9$: tidak ada penyelesaiannya. \quad
$3x^2 = 21$: bagilah dengan 3 lebih dahulu, $x^2 = 7$, penyelesaiannya $\pm\sqrt7$.
\end{example}

\section{Latihan}

\begin{exercise}[$\star$]\label{exo:g9:sqrt:1}
Hitunglah tanpa kalkulator:
\[
\sqrt{64}, \qquad \sqrt{100}, \qquad \sqrt{0.09}, \qquad
\sqrt{\tfrac{1}{25}}, \qquad \left(\sqrt{13}\right)^2, \qquad
\sqrt{(-4)^2} .
\]
\end{exercise}

\begin{exercise}[$\star$]\label{exo:g9:sqrt:2}
Sederhanakan:
\[
\sqrt{50}, \qquad \sqrt{45}, \qquad \sqrt{98}, \qquad \sqrt{300} .
\]
\end{exercise}

\begin{exercise}[$\star$]\label{exo:g9:sqrt:3}
Hitunglah lalu sederhanakan:
\[
\sqrt2 \times \sqrt{18}, \qquad
\frac{\sqrt{75}}{\sqrt3}, \qquad
\sqrt5 \times \sqrt{20}, \qquad
\frac{\sqrt{8}}{\sqrt{50}} .
\]
\end{exercise}

\begin{exercise}[$\star$]\label{exo:g9:sqrt:4}
Selesaikan: $x^2 = 36$; \quad $x^2 = 11$; \quad $x^2 + 4 = 0$; \quad
$2x^2 = 50$.
\end{exercise}

\begin{exercise}[$\star$]\label{exo:g9:sqrt:5}
Sederhanakan menjadi satu suku: $\sqrt{20} + \sqrt{45}$, lalu
$3\sqrt8 - \sqrt{18} + \sqrt2$.
\end{exercise}

\begin{exercise}[$\star\star$]\label{exo:g9:sqrt:6}
Jabarkan lalu sederhanakan:
\[
\left(\sqrt3 + 1\right)^2, \qquad
\left(2\sqrt5 - 3\right)^2, \qquad
\left(\sqrt6 + \sqrt2\right)\left(\sqrt6 - \sqrt2\right).
\]
\end{exercise}

\begin{exercise}[$\star\star$]\label{exo:g9:sqrt:7}
Sebuah ladang berbentuk persegi berluas $6400$ m$^2$. Berapa panjang
sisinya? Berapa diagonalnya (nilai persisnya, lalu \omterm{def:g4:numbers:round}{dibulatkan} ke meter
terdekat)?
\end{exercise}

\begin{exercise}[$\star\star$]\label{exo:g9:sqrt:8}
Tunjukkan bahwa $\dfrac{1}{\sqrt2} = \dfrac{\sqrt2}{2}$ (kalikan
\omterm{def:g4:fractions:def}{pembilang} dan \omterm{def:g4:fractions:def}{penyebutnya} dengan $\sqrt2$). Pakailah trik yang sama
untuk menulis $\dfrac{6}{\sqrt3}$ dan $\dfrac{10}{\sqrt5}$ tanpa \omterm{def:g9:sqrt:def}{akar
kuadrat} pada \omterm{def:g4:fractions:def}{penyebutnya}.
\end{exercise}

\begin{exercise}[$\star\star$]\label{exo:g9:sqrt:9}
Benar atau salah? Berilah alasan dengan bukti atau contoh penyangkal.
\begin{enumerate}
  \item Untuk semua $a, b \geq 0$: $\sqrt{ab} = \sqrt a \sqrt b$.
  \item Untuk semua $a, b \geq 0$: $\sqrt{a+b} = \sqrt a + \sqrt b$.
  \item Untuk semua $x$: $\sqrt{x^2} = x$.
  \item $\left(3\sqrt2\right)^2 = 18$.
\end{enumerate}
\end{exercise}

\begin{exercise}[$\star\star\star$]\label{exo:g9:sqrt:10}
Misalkan $x = \sqrt{7 + 4\sqrt3}$. Hitunglah
$\left(2 + \sqrt3\right)^2$, lalu simpulkan bentuk $x$ yang lebih
sederhana. Pertanyaan yang sama untuk $\sqrt{9 - 4\sqrt5}$ (bidiklah
kuadrat berbentuk $\left(\sqrt5 - 2\right)^2$, dan perhatikan
tandanya).
\end{exercise}

\section{Soal: Bilangan yang bukan pecahan}

\begin{problem}[{Soal akhir pekan --- bukti legendaris bahwa
$\sqrt2$ irasional, banyak sepupunya, dan resep Heron yang luar biasa
baik untuk menghampirinya}]
\label{pb:g9:sqrt:1}
Diagonal persegi satuan adalah panjang yang benar-benar nyata ---
kamu menggambarnya pada \cref{ex:g8:pythagoras:diagonal} --- namun,
seperti yang ditemukan kaum Pythagoras dengan ngeri kira-kira dua
puluh lima abad lalu, \emph{tidak ada \omterm{def:g9:fractions:fraction}{pecahan} apa pun} yang
menakarnya. Legenda menyatakan bahwa penemuan itu dihukum dengan
penenggelaman. Soal ini menuntunmu melewati bukti yang abadi (empat
pertanyaan dan bukti itu menjadi milikmu seumur hidup), memperbanyak
korbannya, lalu berakhir dengan resep yang dipakai para insinyur
selama dua ribu tahun untuk menjinakkan bilangan yang tidak
terjinakkan.

\medskip
\noindent\textbf{Bagian I --- Buktinya.}
\begin{enumerate}
  \item Pertama buruslah bilangannya: periksalah bahwa
        $1.4 < \sqrt2 < 1.5$, lalu bahwa $1.41 < \sqrt2 < 1.42$,
        dengan mengkuadratkan batasnya. Satu angka lagi: di antara
        bilangan tiga desimal yang mana $\sqrt2$ terletak?
  \item Sekarang andaikan --- demi memperoleh kontradiksi --- bahwa
        $\sqrt2 = \frac ab$ untuk suatu \omterm{def:g9:fractions:fraction}{pecahan} yang \emph{paling
        sederhana} (\cref{met:g9:arith:simplify}). Kuadratkan kedua
        ruasnya lalu tunjukkan bahwa $a^2 = 2 b^2$. Bagaimana sifat
        \omterm{def:g3:numbers:evenodd}{ganjil-genap} $a^2$?
  \item Kenyataan \omterm{def:g3:numbers:evenodd}{ganjil-genap} pada \cref{pb:g8:pythagoras:1}
        berkata: \omterm{def:g3:numbers:evenodd}{bilangan ganjil} punya kuadrat yang \omterm{def:g3:numbers:evenodd}{ganjil}.
        Simpulkan bahwa $a$ \omterm{def:g3:numbers:evenodd}{genap}, tulislah
        $a = 2k$, sulihkan --- lalu tunjukkan bahwa $b$ pun pasti
        \omterm{def:g3:numbers:evenodd}{genap}.
  \item Di manakah kontradiksinya? Simpulkan, lalu nyatakan
        teoremanya secara lengkap: \emph{$\sqrt2$ irasional --- ia
        bukan \omterm{def:g3:division:remainder}{hasil bagi} bilangan bulat}.
  \item Buktinya bersandar sekali, secara diam-diam, pada kata
        ``paling sederhana''. Tunjukkan langkah persis mana yang
        akan runtuh tanpa kata itu, lalu jelaskan mengapa
        mengandaikan bentuk paling sederhana itu sah sejak awalnya.
\end{enumerate}

\medskip
\noindent\textbf{Bagian II --- Korbannya berlipat ganda.}
\begin{enumerate}[resume]
  \item Gaya bukti kedua, lewat pemfaktoran prima
        (\cref{thm:g9:arith:factorization}): pada $a^2 = 3b^2$,
        bandingkan sifat \omterm{def:g3:numbers:evenodd}{ganjil-genap} \emph{\omterm{def:g8:powers:def}{eksponen} $3$} pada tiap
        ruasnya (\cref{pb:g9:arith:1}: kuadrat membawa \omterm{def:g8:powers:def}{eksponen} yang
        \omterm{def:g3:numbers:evenodd}{genap}). Simpulkan bahwa $\sqrt3$ irasional.
  \item Jalankan alasan \omterm{def:g8:powers:def}{eksponen} yang sama pada $a^2 = n b^2$ untuk
        bilangan bulat $n$ yang umum: untuk $n$ yang mana alasan itu
        menghasilkan kontradiksi, dan untuk yang mana ia gagal?
        Nyatakan hasil lengkapnya: $\sqrt n$ irasional persis
        ketika\,\dots
  \item Buktikan lema perisainya: bilangan rasional yang bukan \omterm{def:g1:counting:zero}{nol}
        dikali bilangan irasional bersifat irasional. (Andaikan
        $r x = q$ dengan $r, q$ rasional, $r \neq 0$, lalu carilah
        $x$.) Simpulkan bahwa
        $2\sqrt2$ dan $\frac{\sqrt2}{2}$ irasional.
  \item Buktikan bahwa $1 + \sqrt2$ irasional. Lalu yang cantik:
        dengan mengandaikan $s = \sqrt2 + \sqrt3$ rasional,
        hitunglah $s^2$, sendirikan $\sqrt6$, lalu temukan
        kontradiksinya.
  \item Redamlah semangatnya: berikan dua bilangan irasional yang
        \emph{\omterm{def:g1:addition:def}{jumlahnya}} rasional, dan dua yang \emph{\omterm{def:g2:mult:def}{hasil
        kalinya}} rasional. (Keirasionalan tidak terpelihara oleh
        aritmetika --- tiap keadaannya perlu buktinya sendiri.)
\end{enumerate}

\medskip
\noindent\textbf{Bagian III --- Resep Heron.}
Dua ribu tahun sebelum ada kalkulator, Heron dari Aleksandria
menghampiri $\sqrt2$ seperti ini: \emph{tebaklah $x$; gantilah
tebakanmu dengan rata-rata $x$ dan $\frac2x$; ulangi.}
\begin{enumerate}[resume]
  \item Mulailah dari tebakan $x = 1$ lalu hitunglah dua tebakan
        berikutnya sebagai \omterm{def:g9:fractions:fraction}{pecahan} yang persis.
  \item Hitunglah tebakan ketiganya, sekali lagi sebagai \omterm{def:g9:fractions:fraction}{pecahan}
        yang persis, beserta nilai desimalnya. Bandingkan dengan
        pertanyaan 1: berapa desimal $\sqrt2$ yang sudah benar?
  \item Gagasan di balik resepnya: persegi panjang berluas $2$
        dengan satu sisinya $x$ punya sisi yang lain $\frac2x$.
        Tunjukkan bahwa $\sqrt2$ selalu terletak \emph{di antara}
        $x$ dan $\frac2x$ (tinjaulah kedua keadaan $x^2 > 2$ dan
        $x^2 < 2$), jadi merata-ratakan kedua sisinya memeras
        persegi panjangnya menuju sebuah persegi.
  \item Tebakan $\frac32$, $\frac{17}{12}$,
        $\frac{577}{408}$ menyembunyikan permata: hitunglah
        $3^2 - 2 \times 2^2$, lalu $17^2 - 2 \times 12^2$, lalu
        $577^2 - 2 \times 408^2$. Bagian I membuktikan
        $a^2 - 2b^2 = 0$ mustahil --- seberapa dekat \omterm{def:g9:fractions:fraction}{pecahan} Heron
        itu dengan yang mustahil?
  \item Penutupnya, dalam dua kalimat: diagonal persegi satuan itu
        ada di atas kertas, tidak ada \omterm{def:g9:fractions:fraction}{pecahan} yang menakarnya, dan
        penulisan desimalnya tidak pernah dapat berulang
        (\cref{pb:g9:fractions:1}). Jenis ``bilangan baru'' apa yang
        karena itu harus dimuat garis bilangannya --- dan di bagian
        mana pada seri ini kisah lengkapnya diceritakan?

\end{enumerate}
\end{problem}
